%% commands.tex — Eigene Umgebungen und Befehle
%% Alle Farben kommen aus config/colormeta.tex; nichts ist hart kodiert.
%% ──────────────────────────────────────────────────────────────────────────

%% ══════════════════════════════════════════════════════════════════════════
%% 1.  process  +  \step{text}
%%     Nummerierter Schritt-für-Schritt-Ablauf mit farbigem linken Rand.
%%
%%  Verwendung:
%%    \begin{process}
%%      \step{Erster Schritt}
%%      \step{Zweiter Schritt}
%%    \end{process}
%% ══════════════════════════════════════════════════════════════════════════
\newcounter{procstep}

\newenvironment{process}{%
  \setcounter{procstep}{0}%
  \begin{tcolorbox}[
    enhanced,
    breakable,
    boxrule         = 0pt,
    frame hidden,
    borderline west = {3pt}{0pt}{stepcolor},
    colback         = stepcolor!6!white,
    left            = 10pt,
    right           = 8pt,
    top             = 6pt,
    bottom          = 6pt,
    arc             = 3pt,
  ]%
}{%
  \end{tcolorbox}%
}

\newcommand{\step}[1]{%
  \stepcounter{procstep}%
  \vspace{4pt}%
  \noindent
  \colorbox{stepcolor}{%
    \color{white}\bfseries\small\hspace{3pt}\theprocstep\hspace{3pt}%
  }%
  \hspace{6pt}#1\par
}

%% ══════════════════════════════════════════════════════════════════════════
%% 2.  infobox[Titel]
%%     Farbige Infobox mit optionalem fettem Titel.
%%
%%  Verwendung:
%%    \begin{infobox}[Mein Titel] ... \end{infobox}
%%    \begin{infobox}             ... \end{infobox}   % ohne Titel
%% ══════════════════════════════════════════════════════════════════════════
\NewDocumentEnvironment{infobox}{O{}}{%
  \def\lzinfoboxtitle{#1}%
  \ifdefempty{\lzinfoboxtitle}{%
    %% ── Ohne Titel ──
    \begin{tcolorbox}[
      enhanced, breakable,
      colback  = infocolor,
      colframe = infoframe,
      boxrule  = 1pt,
      arc      = 4pt,
      left     = \boxinnersep,
      right    = \boxinnersep,
      top      = \boxinnersep,
      bottom   = \boxinnersep,
    ]%
  }{%
    %% ── Mit Titel ──
    \begin{tcolorbox}[
      enhanced, breakable,
      colback      = infocolor,
      colframe     = infoframe,
      boxrule      = 1pt,
      arc          = 4pt,
      left         = \boxinnersep,
      right        = \boxinnersep,
      top          = \boxinnersep,
      bottom       = \boxinnersep,
      title        = {\textbf{\lzinfoboxtitle}},
      fonttitle    = \bfseries\color{white},
      colbacktitle = infoframe,
      attach boxed title to top left = {yshift=-2mm, xshift=4mm},
      boxed title style = {
        colback  = infoframe,
        colframe = infoframe,
        arc = 3pt, top = 2pt, bottom = 2pt,
      },
    ]%
  }%
}{%
  \end{tcolorbox}%
}

%% ══════════════════════════════════════════════════════════════════════════
%% 3.  warnbox[Titel]
%%     Warnbox / Tipp-Box mit optionalem fettem Titel.
%% ══════════════════════════════════════════════════════════════════════════
\NewDocumentEnvironment{warnbox}{O{}}{%
  \def\lzwarnboxtitle{#1}%
  \ifdefempty{\lzwarnboxtitle}{%
    %% ── Ohne Titel ──
    \begin{tcolorbox}[
      enhanced, breakable,
      colback  = warncolor,
      colframe = warnframe,
      boxrule  = 1pt,
      arc      = 4pt,
      left     = \boxinnersep,
      right    = \boxinnersep,
      top      = \boxinnersep,
      bottom   = \boxinnersep,
    ]%
  }{%
    %% ── Mit Titel ──
    \begin{tcolorbox}[
      enhanced, breakable,
      colback      = warncolor,
      colframe     = warnframe,
      boxrule      = 1pt,
      arc          = 4pt,
      left         = \boxinnersep,
      right        = \boxinnersep,
      top          = \boxinnersep,
      bottom       = \boxinnersep,
      title        = {\textbf{\lzwarnboxtitle}},
      fonttitle    = \bfseries\color{white},
      colbacktitle = warnframe,
      attach boxed title to top left = {yshift=-2mm, xshift=4mm},
      boxed title style = {
        colback  = warnframe,
        colframe = warnframe,
        arc = 3pt, top = 2pt, bottom = 2pt,
      },
    ]%
  }%
}{%
  \end{tcolorbox}%
}

%% ══════════════════════════════════════════════════════════════════════════
%% 4.  procon  +  \pro{text}  \con{text}
%%     Zwei-Spalten-Gegenüberstellung Vorteile / Nachteile.
%%     Alle \pro{} und \con{} werden gesammelt und am Ende nebeneinander
%%     als zwei tcolorbox-Minipages gesetzt.
%%
%%  Verwendung:
%%    \begin{procon}
%%      \pro{Vorteil 1}
%%      \con{Nachteil 1}
%%    \end{procon}
%% ══════════════════════════════════════════════════════════════════════════
\newenvironment{procon}{%
  \def\lzproitems{}%
  \def\lzconitems{}%
  % \def statt \newcommand: kein Fehler beim mehrfachen Aufruf
  % ##1 = #1 auf zweiter Makro-Ebene innerhalb von \newenvironment
  \def\pro##1{\appto\lzproitems{\item ##1}}%
  \def\con##1{\appto\lzconitems{\item ##1}}%
}{%
  \par\smallskip\noindent
  \begin{minipage}[t]{0.475\linewidth}
    \begin{tcolorbox}[
      enhanced,
      colback      = procolor,
      colframe     = proframe,
      boxrule      = 1pt,
      arc          = 4pt,
      title        = {Vorteile},
      fonttitle    = \bfseries\color{white},
      colbacktitle = proframe,
      left = 6pt, right = 6pt, top = 4pt, bottom = 4pt,
    ]
      \begin{itemize}[leftmargin=*, nosep, topsep=0pt]
        \lzproitems
      \end{itemize}
    \end{tcolorbox}
  \end{minipage}%
  \hfill
  \begin{minipage}[t]{0.475\linewidth}
    \begin{tcolorbox}[
      enhanced,
      colback      = concolor,
      colframe     = conframe,
      boxrule      = 1pt,
      arc          = 4pt,
      title        = {Nachteile},
      fonttitle    = \bfseries\color{white},
      colbacktitle = conframe,
      left = 6pt, right = 6pt, top = 4pt, bottom = 4pt,
    ]
      \begin{itemize}[leftmargin=*, nosep, topsep=0pt]
        \lzconitems
      \end{itemize}
    \end{tcolorbox}
  \end{minipage}%
  \par\smallskip
}

%% ══════════════════════════════════════════════════════════════════════════
%% 5.  codebox
%%     Monospace-Block mit grauem Hintergrund (für kurze Ausschnitte).
%%     Für Syntax-Highlighting: listings-Paket nutzen (siehe packages.tex).
%% ══════════════════════════════════════════════════════════════════════════
\newenvironment{codebox}{%
  \begin{tcolorbox}[
    enhanced,
    breakable,
    colback  = codeblock,
    colframe = codeframe,
    boxrule  = 0.5pt,
    arc      = 3pt,
    left     = 8pt,
    right    = 8pt,
    top      = 6pt,
    bottom   = 6pt,
    fontupper = \ttfamily\small,
  ]%
}{%
  \end{tcolorbox}%
}

%% ══════════════════════════════════════════════════════════════════════════
%% 6.  \sidetables{Inhalt links}{Inhalt rechts}
%%     Zwei beliebige Inhalte (z. B. Tabellen, Listen) nebeneinander.
%% ══════════════════════════════════════════════════════════════════════════
\newcommand{\sidetables}[2]{%
  \par\smallskip\noindent
  \begin{minipage}[t]{0.475\linewidth}
    #1
  \end{minipage}%
  \hfill
  \begin{minipage}[t]{0.475\linewidth}
    #2
  \end{minipage}%
  \par\smallskip
}

%% ══════════════════════════════════════════════════════════════════════════
%% 7.  stdtable{Spaltendefinition}{Kopfzeile}
%%     Formatierte Tabelle mit farbiger Kopfzeile und alternierenden Zeilen.
%%     Spaltentypen: l r c  für Text; X für auto-breite Spalten; p{…} für
%%     feste Breite. X-Spalten sind durch tabularx möglich, weil \NewEnviron
%%     (environ-Paket) den Body komplett sammelt, bevor tabularx ihn scannt.
%%
%%  Verwendung:
%%    \begin{stdtable}{l X X}{Algorithmus & Kompl. & Notiz}
%%      Bubble Sort & $O(n^2)$ & Einfach \\
%%    \end{stdtable}
%% ══════════════════════════════════════════════════════════════════════════
\NewEnviron{stdtable}[2]{%
  \par\smallskip
  \rowcolors{2}{tablerowalt}{white}%
  \begin{center}%
  \begin{tabularx}{\linewidth}{#1}%
  \rowcolor{tableheader}%
  \color{tableheadertext}\bfseries #2\\[\jot]%
  \BODY
  \end{tabularx}%
  \end{center}%
  \rowcolors{1}{}{}%   Reset: verhindert Übertrag auf nachfolgende Tabellen
  \par\smallskip
}
